\begin{ejer}[Reality check]
$\mathbb{Z}_{(2)}$ y $\mathbb{Z}_2$ (los enteros $2$-ádicos) son ambos anillos de valuación discreta con cuerpo residual $\mathbb{F}_2$, y sus espectros tienen la misma forma. ¿Son el mismo anillo? Da un elemento de $\mathbb{Z}_2$ que no esté en $\mathbb{Z}_{(2)}$, y explica por qué $\mathbb{Z}_{(2)}\subset\mathbb{Z}_2$ pero $\mathbb{Z}_{(2)}\neq\mathbb{Z}_2$.
\end{ejer}
\begin{sol}
$\mathbb{Z}_2$ se define como límite inverso de
\[
\begin{tikzcd}[column sep=large]
\cdots \arrow[r,"\rho_{4,3}"]
    & \mathbb{Z}/2^3\mathbb{Z} \arrow[r,"\rho_{3,2}"]
    & \mathbb{Z}/2^2\mathbb{Z} \arrow[r,"\rho_{2,1}"]
    & \mathbb{Z}/2\mathbb{Z}.
\end{tikzcd}
\]
Esto significa que
\[
\mathbb{Z}_2=\set*{(z_n)_{n\en\mathbb{N}}\given
z_n\en\mathbb{Z}/2^n\mathbb{Z}\;\text{ y }\;\rho_{n+1,n}(z_{n+1})=z_n}
\]
Para cada $n\en\mathbb{N}$, el mapa cociente $q_n\colon\mathbb{Z}\to\mathbb{Z}/2^n\mathbb{Z}$ conmuta con los $\rho_{m,n}$, ya que estos se derivan así a partir del cociente. Se sigue de la propiedad universal del límite que los $q_n$ inducen un único morfismo $q\colon\mathbb{Z}\to\mathbb{Z}_2$ con $\pi_n\circ q=q_n$. De forma concreta,
\[
q(z)=(q_n(z))_{n\en\mathbb{N}}
\qquad\text{para todo }z\en\mathbb{Z}.
\]
Notemos que $q$ es mono, pues si $q(z)=0$ se sigue que $q_n(z)=0$ para todo $n\en\mathbb{N}$, así que $z\equiv0\pmod{2^n}$ para todo $n\en\mathbb{N}$. La única forma de que esto ocurra es que $z=0$.
\vspace{1em}\\
\emph{Veamos que $q\colon\mathbb{Z}\to\mathbb{Z}_2$ induce un
$\widetilde q\colon\mathbb{Z}_{(2)}\to\mathbb{Z}_2$}:\\ 
Basta verificar que
$q(z)$ es invertible para $z\noten(2)$, i.e., para $z$ impar. En efecto,
como $\gcd(z,2^n)=1$, $q_n(z)\en\mathbb{Z}/2^n\mathbb{Z}$ es invertible y
\[
\rho_{n+1,n}\bigl(q_{n+1}(z)\bigr)=q_n(z).
\]
Como los $\rho_{n+1,n}$ son homomorfismos de anillos, se sigue que
\[
\rho_{n+1,n}\bigl(q_{n+1}(z)^{-1}\bigr)=q_n(z)^{-1}
\]
para todo $n\en\mathbb{N}$. Esto prueba que
$(q_n(z)^{-1})_{n\en\mathbb{N}}\en\mathbb{Z}_2$, siendo este el inverso de
$q(z)$.

Finalmente, gracias a que $q(z)$ es invertible para $z\noten(2)$, $q$ induce un $\widetilde q\colon\mathbb{Z}_{(2)}\to\mathbb{Z}_2$; más aún, como $q$ es mono, también lo es $\widetilde q$.
\vspace{1em}\\
\emph{Construyamos la expansión binaria de $\mathbb{Z}_2$}:\\ 
Para todo $(c_n)_n\en\mathbb{Z}_2$, sea $\widehat c_n\en\set{0,1,\ldots,2^n-1}$ el
representante canónico de $c_n\en\mathbb{Z}/2^n\mathbb{Z}$, y convengamos $\widehat c_0=0$. La condición
$\rho_{n+1,n}(c_{n+1})=c_n$ implica
\[
c_{n+1} \in \rho^{-1}_{n+1,n}(\set{c_n}) = \set[\big]{[\widehat c_n]_{2^{n+1}},[\widehat c_n + 2^n]_{2^{n+1}}},
\quad\text{es decir,}\quad
\widehat c_{n+1}\en\set{\widehat c_n,\widehat c_n+2^n}
\]
Esto muestra que un elemento de $\mathbb{Z}_2$ tiene solo dos elecciones
posibles para cada entrada al haber escogido las anteriores. Así, hay una
biyección
\[
\xi\colon\mathbb{Z}_2\longrightarrow\set{0,1}^{\mathbb{N}},
\qquad
\xi((c_n)_n)=\left(2^{-n}(\widehat c_{n+1}-\widehat c_n)\right)_{n=0}^{\infty}.
\]
(Si $(x_n)_n=\xi(c)$, esta sucesión se denota como la cadena infinita de dígitos binarios $\cdots x_nx_{n-1}\cdots x_0$).

Notemos que la inversa está dada por
\[
(x_n)_n\longmapsto
\left(\left[\sum_{i=0}^{n-1}x_i2^i\right]_{2^n}\right)_n
\]
En efecto,
\[
\rho_{n+1,n}\left(\left[ \sum_{i=0}^{n}x_i2^i \right]_{2^{n}}\right)
=\left[ x_n2^n+\sum_{i=0}^{n-1}x_i2^i \right]_{2^{n+1}}
=\left[ \sum_{i=0}^{n-1}x_i2^i \right]_{2^{n+1}}
\]
de modo que $\left(\left[\sum_{i=0}^{n-1}x_i2^i\right]_{2^n}\right)_{n\en\mathbb{N}}\in\mathbb{Z}_2$; mientras que sus composiciones son en cada entrada:
\[
\left[ \sum_{i=0}^{n-1}\left(2^{-i}(\widehat c_{i+1}-\widehat c_i)\right)2^i \right]_{2^n}
= \left[ \sum_{i=0}^{n-1}(\widehat c_{i+1}-\widehat c_i) \right]_{2^{n}} 
= \left[ \widehat c_n \right]_{2^{n}} = c_n
\]
y
\[
2^{-n}\left( \hat{\left[\sum_{i=0}^{n}x_i2^i\right]}_{2^{n}} - \hat{\left[\sum_{i=0}^{n-1}x_i2^i\right]}_{2^{n}}\right)
= 2^{-n}\left( \sum_{i=0}^{n}x_i2^i - \sum_{i=0}^{n-1}x_i2^i\right)
=2^{-n}(x_n2^n)=x_n
\]
Esto prueba que son inversos entre sí. (Por esta inversa es que al número con dígitos $\cdots x_nx_{n-1}\cdots x_0$ se le denota $\sum_{i=0}^{\infty}x_i2^i$; la convergencia de esta serie formal se justifica mediante la norma $2$-ádica).
\vspace{1em}\\ 
\emph{Todo elemento de $\widetilde q(\mathbb{Z}_{(2)})$ tiene expresión binaria eventualmente periódica}:
\begin{itemize}
    \item Si $a\en\mathbb{N}_0$ tiene expresión binaria finita $a=\sum_{i=0}^{N}a_i2^i$, $a_i\en\set{0,1}$, al poner $a_i=0$ para $i>N$:
    \[
    q_n(a)=\left[\sum_{i=0}^{n-1}a_i2^i\right]_{2^{n}}
    \quad\Longrightarrow\quad
    \xi(q(a))=\cdots 0a_Na_{N-1}\cdots a_0,
    \]
    periódica a partir de $i>N$ con $a_i=0$.

    \item Sean $a,b\en\mathbb{N}_0$ con $b$ impar y $\gcd(a,b)=1$. Como $\gcd(2,b)=1$, el teorema de Euler--Fermat implica que existe $k\en\mathbb{N}$ con $2^k\equiv1\pmod b$; esto es, $b\mid(2^k-1)$. Luego existe $m\en\mathbb{N}$ con $2^k-1=bm$. Ahora bien,
    \[
    \left(\sum_{i=0}^{n-1}2^{ik}\right)(1-2^k)
    =\sum_{i=0}^{n-1}2^{ik}-2^{(i+1)k}
    =1-2^{nk}
    =1\pmod{2^{nk}}.
    \]
    Denotando $\widetilde q_j\coloneq\pi_j\circ\widetilde q$ (es decir, $\tilde q=(\tilde q_j)_j$), esto prueba que para todo $n\en\mathbb{N}$:
    \[
    \widetilde q_{nk}\left(\frac{1}{1-2^k}\right)=\left[\sum_{i=0}^{n-1}2^{ik}\right]_{2^{nk}}
    \]
    Más aún, como para $0<t\leq k$,
    \begin{align*}
    \widetilde q_{nk-t}\left(\frac{1}{1-2^k}\right)
    =\rho_{nk,nk-t}\left(\widetilde q_{nk}\left(\frac{1}{1-2^k}\right)\right)
    =\rho_{nk,nk-t}\left(\left[\sum_{i=0}^{n-1}2^{ik}\right]_{2^{nk}}\right)
    =\left[\sum_{i=0}^{n-1}2^{ik}\right]_{2^{nk-t}}
    \end{align*}
    tenemos que
    \[
    \widetilde q_\ell\left(\frac{1}{1-2^k}\right)
    =\left[\sum_{i=0}^{n-1}2^{ik}\right]_{2^{\ell}}
    \qquad\text{para }(n-1)k\leq \ell<nk.
    \]
    por lo cual $\xi\left(\widetilde q\left(\frac{1}{1-2^k}\right)\right) = \cdots10\cdots01\underbrace{0\cdots0}_{\mathclap{k-1\;\text{ceros}}}1 = \ubar{0\cdots01}$ es eventualmente periódica.

    Por último, se puede probar, analizando acarreos al sumar bloques periódicos, que la suma de $2$-ádicos eventualmente periódicos es eventualmente periódica, así que
    \[
    \xi\left(\widetilde q\left(-\frac{a}{b}\right)\right)
    =\xi\left(\widetilde q\left(\frac{am}{1-2^k}\right)\right)
    =\xi\left(\sum_{s=1}^{am}\widetilde q\left(\frac{1}{1-2^k}\right)\right)
    \]
    es eventualmente periódica. Esto prueba que, para todo $z\en\mathbb{Z}_{(2)}$ no positivo, $\xi(\widetilde q(z))$ es eventualmente periódico.

    \item Sea $z\en\mathbb{Z}_{(2)}$ positivo. Como $-z-1<0$, $\xi(\widetilde q(-z-1))=(a_i)_{i=0}^{\infty}$ es eventualmente periódico, así que al definir $(a_i')_{i=0}^{\infty}\en\set{0,1}^{\mathbb{N}}$ mediante $a_i'+a_i=1$ para todo $i\en\mathbb{N}$:
    \[
    q_n(-z-1)+ \left[\sum_{i=0}^{n-1}a_i'2^i\right]_{2^{n}}
    = \left[\sum_{i=0}^{n-1}(a_i+a_i')2^i\right]_{2^{n}}
    = \left[\sum_{i=0}^{n-1}2^i\right]_{2^{n}}
    \]
    implica
    \[
    \bigl( \underbrace{q_n(-z-1)+1}_{= -q_n(z) } \bigr)  +  \left[ \sum_{i=0}^{n-1}a_i'2^i \right]_{2^{n}}
    = \left[\sum_{i=0}^{n-1}2^i+1\right]_{2^{n}}
    =\left[2^n\right]_{2^{n}}
    =\left[0\right]_{2^{n}}
    \]
    por unicidad de inversos aditivos, esto prueba que $q_n(z)=\sum_{i=0}^{n-1}a_i'2^i$ para todo $n\en\mathbb{N}$, por lo que $\xi(\widetilde q(z))=(a_i')_{i=0}^{\infty}$ es eventualmente periódico.
\end{itemize}

En conclusión, como todo elemento de $\tilde q(\mathbb{Z}_{(2)})$ tiene expansión binaria
eventualmente periódica, basta con definir uno de $\mathbb{Z}_2$ que no la tenga. Por ejemplo, si definimos $(x_i)^\infty_{i=1} \en \set{0,1}^{\mathbb{N}_0}$ por
\[
x_i \coloneq 
\begin{cases}
    1, &\text{si existe $k\en\mathbb{N}_0$ con $i^2=k$}\\
    0, &\text{en otro caso.}
\end{cases}
\]
Esto es, $(x_i)_i$ tiene un $1$ exactamente en las posiciones que son
cuadrados perfectos. Como las distancias entre cuadrados consecutivos crecen
sin cota, esta expansión no es eventualmente periódica. Por tanto,
$\xi((x_i)_i)\in \mathbb{Z}_2\setminus\widetilde q(\mathbb{Z}_{(2)})$, lo cual garantiza que
$\mathbb{Z}_{(2)}\hookrightarrow\mathbb{Z}_2$ no es sobre.
\end{sol}